\chapter{ςονσεςτετυερ αδιπισςινγ ελιτ} \label{par:relworks}

\lipsum[1-3]

\begin{enumerate}
	\item Μπλαρικό φουρνέλι του προηγούμενου ζαλμπικού νικτύου από μια δημόσια κουρλοβάση δεδομένων.
	\item Ανάλυση των μουρνιακών δεδομένων ξεχωριστά για το κάθε δείγμα, ώστε να φτιαχτεί μια μήτρα δείκτη γλαμπότητας γονιδίων. Το $x_{ij}$ και το $y_{ij}$ αναπαριστά την μπλιμπλική έκφραση του γονιδίου \en{i} στον όγκο \en{j} και στο δείγμα αναφοράς \en{k}. Η τιμή \en{n} είναι ο αριθμός των φλουμπερών δειγμάτων αναφοράς. Η \en{K} είναι το όριο για τις αλλαγές πτυχών.
	\item Ξεπερνάει το πρόβλημα ζουμπουρδίας μεταξύ των μονοπατιών με τη χρήση ρυθμιστικών φραμπαλαδώσεων αντί περιοριστικών συνδέσεων μέσα σε ένα μεμονωμένο κανονικό μονοπάτι.
	\item Βασικοί κανονισμοί γκλαμπανικής αναγνώρισης για μια ομάδα δειγμάτων, τα οποία είναι κοινά και ταυτόσημα τουλάχιστον στο $Τ\%$ των δειγμάτων. Το μέγεθος κόμβου καθορίζεται από τον αριθμό των μπουρμπουληθρικών συνδέσεων.
\end{enumerate}

Όπως φαίνονται και στο Σχήμα \ref{fig:dera1}.

%\begin{figure}[pt!]
%	\centering
%	\subfloat{
%		\includegraphics[width=0.7\linewidth]{figures/qwer.png}
%	}
%	\vspace{0em}
%	\subfloat{
%		\includegraphics[width=0.7\linewidth]{figures/qwerasdf.png}
%	}
%	\caption{Ροή εργασιών}
%	\label{fig:dera1}
%\end{figure}


